\documentclass[10pt,letterpaper]{article}
\usepackage[top=1in,left=1in,right=1in,bottom=1in]{geometry}
\usepackage{microtype}
\usepackage{graphicx}
\usepackage{booktabs}
\usepackage{xcolor}
\usepackage[round,authoryear]{natbib}
\usepackage[colorlinks=true,allcolors=blue]{hyperref}

\setlength{\parindent}{0pt}
\setlength{\parskip}{1em}

\title{\vspace*{-0.8in}\textbf{Response to Reviewers --- Second Revision}}
\author{}
\date{}

\begin{document}
\maketitle

\noindent\textbf{Manuscript:} A Four-Dimension Pre-Modeling Audit Protocol for Educational Prediction Benchmarks\\
\textbf{Journal:} Scientific Reports\\
\textbf{Revision:} Second round

\vspace{0.2in}
\noindent Dear Editor and Reviewers,

We thank the reviewers for their careful reading of our revised manuscript. Below we address each remaining comment point by point. All changes are described with specific line references.

\section*{Response to Reviewer 1}

\textbf{Recommendation:} Accept with minor suggestion.

\textbf{Comment:} The authors may wish to cite Saşmaz (2026), which uses a multi-model machine learning framework on PISA 2022 data and is relevant to the benchmarking methodology discussed here.

\textbf{Response:} We thank the reviewer for this suggestion. We have added a critical citation of \citet{Sasmaz2026} to the Discussion section (line~351). The added sentence notes that, while \citet{Sasmaz2026} applies a multi-model framework with stratified $k$-fold cross-validation to PISA 2022 data, it does not report group-aware holdout results, illustrating the persistent gap between current evaluation practice and the audit standards proposed here. We believe this framing strengthens rather than dilutes our argument.

\section*{Response to Reviewer 2}

\textbf{Recommendation:} Minor revision (second round).

\vspace{0.2in}
\noindent We thank Reviewer~2 for the close reading and for pushing us to clarify four important points.

\subsection*{Comment 2-1: Feature set alignment for fair retention-ratio calculation}

\textbf{Reviewer:} The retention ratio currently compares group-holdout $R^2$ (with group identifiers excluded) against iid $R^2$ (with group identifiers included). This asymmetry needs to be addressed.

\textbf{Response:} We agree and have made the following changes:

\begin{itemize}
  \item Modified \texttt{run\_7dataset\_audit.py} (lines~238--244) to run IID splits on the same clean feature matrix (group-identifier columns excluded) when the adapter reports such columns. The clean iid $R^2$ is stored as \texttt{iid\_r2\_clean} in the output profile.
  \item Updated the retention-ratio logic (line~337) to use \texttt{iid\_r2\_clean} when available, ensuring apples-to-apples comparison.
  \item Re-ran the full pipeline and updated all affected numbers:
    \begin{itemize}
      \item xAPI-Edu: iid $R^2$ 0.599 $\to$ 0.614, retention 81\% $\to$ 79\% (profile unchanged: Strong)
      \item UCI Student: iid $R^2$ 0.262 $\to$ 0.242, retention $-37\%$ $\to$ $-40\%$ (profile unchanged: Mostly Passing)
      \item Entrance Exam: iid $R^2$ unchanged at 0.438 (board features have negligible impact on iid performance)
    \end{itemize}
  \item Updated all affected numbers in the main text (xAPI-Edu 0.614, UCI Student 0.242, retention ratios 79\% and $-40\%$).
  \item Removed stale values (0.599, 0.262, 81\%, $-37\%$) from the manuscript.
  \item Updated \texttt{results/audit\_7dataset\_results.json} with the corrected canonical values.
  \item Regenerated all figures.
\end{itemize}

No profile boundaries were crossed for any dataset, and the three key substantive conclusions remain unchanged.

\subsection*{Comment 2-2: Structural versus model-dependent fragility}

\textbf{Reviewer:} Ridge and ensemble models recover group-holdout performance for Entrance Exam. Failure of the unregularized linear model alone is insufficient to establish structural fragility. Persistent failures across model families should be distinguished from model-dependent or partly remediable failures.

\textbf{Response:} We fully agree and have revised the manuscript throughout to draw this distinction clearly:

\begin{itemize}
  \item \textbf{Finding~1} (line~333): Now explicitly distinguishes ``structural fragility'' (Higher Ed: all models $R^2 < -8$; UCI Student: linear $-0.097$, RF $0.018$, GBT $-0.057$) from ``linear-model-only failure that is partly remediable by ensemble models'' (Entrance Exam: RF $R^2 \approx 0.32$ vs. linear $-0.060$).
  \item \textbf{Finding~2} (line~335): Retitled to ``Structural fragility persists across model families, while model-dependent failures are partly remediable.'' Includes Entrance Exam's ensemble recovery as a counterexample.
  \item \textbf{Finding~4} (line~339): Now notes that Entrance Exam's iid-to-group gap is ``partly closed by ensemble models.''
  \item \textbf{Conclusion~(1)} (line~369): Separates structural collapse from model-dependent failure.
  \item \textbf{Conclusion~(2)} (line~371): Retitled from ``Fragility persists across model families'' to ``Structural collapse versus model-dependent failure.''
  \item \textbf{Detailed-results section} (line~194): Adds sentence: ``The nature of the failure differs: UCI Student and Higher Ed exhibit structural fragility (all model families collapse), whereas Entrance Exam's failure is linear-model-specific and partly remediable---its RF group $R^2$ reaches approximately $0.32$.''
  \item \textbf{Table~1 caption} (line~188): Replaced ``flags the structural fragility of the base benchmark'' with language describing the audit's role as flagging the weakest link, with the detailed breakdown distinguishing structurally invariant collapse from model-dependent failure.
  \item \textbf{Line~260}: Replaced ``fragile or partially fragile datasets'' with ``datasets with metadata-adequacy failures'' and annotated each dataset's failure pattern.
  \item Remaining uses of ``fragility'' in the abstract and general discussion are retained as high-level descriptors of the cross-group generalization failure phenomenon.
\end{itemize}

\subsection*{Comment 2-3: Feature-attribution scope limitation}

\textbf{Reviewer:} The feature-attribution analysis is computed under iid splits only. The connections drawn to audit dimensions are interesting but not directly supported. Either compute group-aware attributions or clearly limit scope.

\textbf{Response:} We have chosen to clearly limit the scope rather than write new group-aware attribution code. At the end of the feature-attribution section (line~322), we added:

\emph{``We note that this feature-attribution analysis is conducted under iid random splits and does not directly evaluate attribution stability under group-aware holdout; the connections drawn to the audit profiles describe the alignment between split-level attribution stability and dataset readiness rather than establishing a causal link to cross-group behavior.''}

\subsection*{Comment 2-4: RF importance type and feature-name mapping}

\textbf{Reviewer:} (a) Report the type of feature importance used (impurity-based or permutation-based). (b) Add a table mapping feature codes to meaningful names.

\textbf{Response:} (a) In the feature-attribution section (line~318), we now specify: ``Random Forest feature importances (impurity-based, using scikit-learn's default \texttt{feature\_importances\_}).'' 

(b) We added a new Supplementary Table~S7 (page~9 of the supplementary PDF) mapping all feature codes (e.g., f5, f28, f3) used in Tables~S5--S6 to human-readable feature names for UCI Student, Higher Ed, and OULAD. The table covers every feature that appears in the top-10 coefficient and importance tables.

\section*{Response to Reviewer 3}

\textbf{Recommendation:} Not assigned to this revision.

\textbf{Response:} As noted in our previous response, the scope and methods of this paper---a pre-modeling audit protocol for educational prediction datasets---remain distinct from automated short-answer scoring or LLM-based assessment. We respectfully reiterate that Reviewer~3's concerns, which focus on LLM grading reliability and rubric design, address a different task and a different stage of the educational AI pipeline. We defer to the editor's judgment on whether the current reviewer panel is appropriate for this manuscript.

\vspace{0.3in}
\noindent We hope these revisions fully address all remaining concerns. We are grateful for the reviewers' constructive feedback, which has substantially improved the clarity and precision of the manuscript.

\vspace{0.1in}
\noindent Respectfully,\\
Yan Ma \and Lizhuo Zhang

\newpage
\bibliographystyle{unsrtnat}
\bibliography{behavioraudit}

\end{document}
